\chapter{The Native Apparatus Presents The Residue}

\section{Forcing the residue to show}

The last three chapters split the number into faces and turned it on a plane. This part turns to the residue's
reckoning --- to what the machine finally \emph{does} with the leftover it has carried since the beginning. Recall what
it cannot do. It cannot discard the residue: the step down to nothing was proved false, an uninhabited case the checker
will not compute. And it cannot store the residue whole: the machine keeps finite conditions, never a completed
infinity. So a third thing is forced, and this section is about that third thing. The apparatus \emph{presents} the
residue --- forces it into view --- and it does so not by any special act but by the plainest operation it performs
every time it reads.

Start with how the native apparatus reads at all, because everything follows from its being a ratio. The apparatus
does not read by naming a magnitude; it reads by \emph{dividing} --- by comparing one count against another and
reporting the fraction. Its reading is a numerator over a denominator, a ratio of two counts. And the denominator is
worth a word, because it is the apparatus's \emph{own} scale, not a unit borrowed from the world: the machine divides
by a count it set for itself, and its reading is honest against that count, not against any external ruler. To read,
for this apparatus, is simply to take a ratio.

Now what a ratio yields, which is two things and not one. Divide one count by another and you get a whole part --- how
many times the denominator goes in cleanly --- and a leftover: what is left when the whole part is taken out, the piece
too small to make one more of the divisor. The machine names the first the \emph{floor} and the second the
\emph{remainder}, and it computes both. The floor is the whole the apparatus keeps --- the clean part of the reading,
the count that divided out. The remainder is the part the division could \emph{not} resolve into a whole: the fraction
left below the scale, the piece the ratio could not absorb. Every reading the apparatus makes is a floor plus a
remainder, and the remainder is exactly the residue --- the leftover the book has tracked from its first chapter ---
now written in the apparatus's own arithmetic.

And here is the forcing, which is why the section has its name. The apparatus, having a remainder, has only honest
options. It cannot throw the remainder away: the earlier chapter proved that the step down to nothing does not
typecheck. It cannot round the ratio and call it whole: that would be to report an exact reading where the division
left a leftover, which is a lie about being exact. So the only honest thing left is to \emph{present} the remainder ---
to compute the leftover and show it. ``The native apparatus presents the residue'' is not a figure of speech; it is a
description of the arithmetic. The remainder \emph{is} the residue, and it is surfaced by the very division the
apparatus performs in order to read at all. The floor is the whole the machine keeps; the remainder is the residue it
is \emph{forced} to show --- forced not by a rule imposed from outside but by the plain fact that a ratio which hides
its remainder is a ratio lying about being exact. The residue shows because the division that reads it cannot help but
leave it.

This is the general act, and the rest of the chapter watches it happen on two particular benches. On one, the
apparatus is a physical balance, and the residue it presents is read, by the physical gauge, as a measured quantity;
that bench and its reading belong to the next volume. On the other, the apparatus is the checker itself, and the
residue is presented under elaboration; that bench belongs to the computational gauge. This section keeps only the
arithmetic beneath both: whatever the bench, the apparatus's remainder is the residue, forced to show by the division
that reads.

In the two verbs the split is exact. The floor is a \emph{funge}: the whole part, everything that bagged cleanly into a
single quotient, gathered by the division into one clean count. The remainder is a \emph{tange}: the residue, selected
out by the same division as precisely the part that would \emph{not} bag --- the fraction that could not be gathered
into the whole. Forcing the residue to show is the operation that tanges the leftover out of the whole and cannot put
it back. The apparatus cannot funge the remainder into the floor; the division that produced the clean count produced,
in the same act, the piece that would not join it. So the residue is tanged into view by the very reading that funges
the rest away.

And this is where the ending first comes into focus. The residue the apparatus is forced to show is the same residue
the whole instrument will, at its far end, ask for a number. What the closing chapters add is only which reading, on
which bench, at what floor. And the honesty the book ends on is already visible in this small operation: a machine that
presents its remainder rather than rounding it away is a machine that will, at the last, hand back a bracket rather
than force a point. The open interval the book closes on is this same remainder, held rather than hidden --- the
leftover the apparatus was always forced to show, carried to the last reading the instrument makes.

\section{The physical bench}

The last section gave the apparatus in general: a ratio, a floor, and a remainder forced to show. This short section
watches that apparatus become a \emph{physical} balance --- and it is deliberately brief, because most of what happens
on the physical bench belongs to the next volume, not this one. What this volume keeps is small, and it is worth
keeping.

The first thing to keep is that the physical bench is the \emph{same} apparatus. When the machine is set up as a
balance --- weighing one mass against another --- its reading is still a ratio, still a floor plus a remainder, and the
residue it presents is still the remainder. Nothing new is built to make it physical; the balance is the ratio of the
last section, pointed at masses. And there is a small, exact reward in the setup: the mass the balance weighs against
--- its test mass --- is the electron the machine named two chapters ago. The thing forced into its box by the
pigeonhole, named by counting, is the very thing the physical bench puts on its pan. The named object and the weighed
object are one.

The second thing to keep is the most important, and it is the blind discipline holding firm at the one place you might
expect it to break. A physical balance, you would think, finally reaches outside the machine --- reaches a real force,
a real constant, a number in the world. It does not. The construction is explicit that what the physical bench performs
is a \emph{calibration of the device's own reading}, not a measurement in external units: the masses, the reading, the
whole setup are made visible through the same internal wheel the machine has used all along --- its own scale of
charge and mass and value. Even here, weighing real masses on a real balance, the machine calibrates itself against its
own floor, through its own wheel. It does not reach an external constant. There is no outside number on this bench
either. And that is the deepest reason the book's ending can be honest: the machine never had an external value to hand
you --- not on the computational bench, and not, it turns out, even on the physical one. Every bench measures the
machine to its own floor.

The third thing is a boundary, not a content. The physical reading itself --- the balance as a torsion instrument, the
source masses, the constant of gravitation the setup is arranged to calibrate --- is the physical gauge's to interpret,
and it belongs to the next volume. This volume names the bench and stops there. It is a Cavendish balance; what its
residue \emph{means}, read as a physical quantity, is Volume Two's cut, not this one's. The point of naming it here is
only to show that the same residue-presenting apparatus wears a physical face --- and that even wearing it, it stays
inside the machine's own scale.

In the two verbs the physical bench adds nothing new and confirms what the last section said. The whole physical setup
--- masses, balance, wheel --- is funged into the one scale the machine reads on; and the residue is still the
remainder, tanged into view by the division, exactly as before. The bench is physical; the arithmetic is not. So Part
IV's opening lesson survives its hardest test: pointed at real masses, the apparatus still presents a remainder against
its own floor, and still reaches no number outside itself. The next section turns the same apparatus toward
computation, where the residue shows once more --- in a different currency, on the same honest terms.

\section{The computational bench}

The last section watched the apparatus become a physical balance and present its residue as a remainder against its own
scale. This closing section turns the same apparatus a second way --- toward computation --- and finds the residue
shown once more, in a currency we have already met. It is a short section, like the last, because most of what it opens
belongs to another volume. But it closes the chapter on the point the whole part has been building: the residue shows
on \emph{every} bench, and the benches agree.

The computational bench is, again, the \emph{same} apparatus --- the same ratio, the same floor, the same remainder
forced to show. What changes is only the currency the reading is billed in. On this bench the machine is the checker
itself, and to present a term is to \emph{elaborate} it --- to run the reduction that settles it --- and elaboration,
as an earlier chapter established, costs the machine beats of its own clock. So the residue shows here as a \emph{cost}:
what the elaboration spends and cannot recover, the remainder left over when the settled part of the work is
subtracted, counted not in the divisions of a balance but in the machine's own heartbeats. The apparatus divides here
too; it just divides time.

And now the point Part IV was arranged to reach, stated with the care it needs. The residue has now shown on two
benches: a physical balance in the last section, and the checker's own elaboration in this one. These two benches share
no vocabulary. One weighs masses on a torsion fiber; the other counts heartbeats of computation; neither can see what
the other does, and neither was built from the other. And yet the leftover they present is, the book reads, the
\emph{same} residue --- the one un-eliminable term, surfacing on a balance and surfacing in a cost. That agreement is
the evidence the whole instrument is built to earn: two decorrelated readings of one object, landing together. But the
honesty here has to be exact, and it is the honesty the earlier chapter on the faces already established. The machine
does not \emph{prove} a cross-bench identity; there is no theorem that the balance's remainder equals the elaborator's.
What each bench proves is its own remainder; that the two are the same residue is the book's reading of the situation
--- consilience, the agreement of independent instruments --- and not a computation the machine performs. It is exactly
as strong as agreement across decorrelated readings ever is, and no stronger; and that is precisely why it is worth
trusting: it cannot have been arranged.

What the elaboration residue \emph{means}, read as a computation --- a cost of proof, a price of settling --- is the
computational gauge's cut, and it belongs to the third volume; this one names the bench and hands the reading forward.
In the two verbs the bench adds its confirmation: elaboration funges the settled work into a count of beats, and the
residue is the remainder that count leaves over, tanged into view exactly as the balance's remainder was. Two benches,
two currencies, one leftover --- funged differently, tanging out the same residue.

So Part IV closes where it aimed. The residue could not be discarded and could not be stored whole, so the apparatus
was forced to present it --- and it presents it on whatever bench you build, physical or computational, against its own
floor, reaching no number outside itself. That two such different benches agree on the leftover is the reason the
reading can be trusted, and it is the reason the book's ending, when it holds its bracket open, will be believed: not
because one gauge insists, but because independent instruments, sharing nothing, keep finding the same residue and
keep, honestly, refusing to round it away.
